\documentclass[10pt,a4paper]{amsart}
\usepackage[margin=19mm]{geometry}
\usepackage{amsmath,amssymb,mathtools}
\usepackage[scr=rsfs,cal=euler]{mathalfa}
\usepackage{xfrac}
\usepackage[unicode,hidelinks,bookmarks=false]{hyperref}
\newcommand{\ie}{i.e.\ }
\newcommand{\cf}{cf.\ }
\newcommand{\resp}{respectively\ }
\newcommand{\Subsection}{\subsection}
\providecommand{\nobreakdash}{\nobreak\hbox{-}}
\setlength{\emergencystretch}{2em}
\multlinegap0pt
\usepackage{mathtools}
\usepackage{amssymb}
\usepackage{bm}
\usepackage{enumerate}
\usepackage{xcolor}
\usepackage{float}
\newtheorem{lemma}{Lemma}
\newtheorem{proposition}{Proposition}
\newcommand{\depth}{\operatorname{depth}}
\newcommand{\m}{\mathfrak{m}}
\newcommand{\wt}{\widetilde}
\title{Integrally closed ideals with $e_{2}(I)=e_{1}(I)-e_{0}(I)+\lambda(A/I)$}
\author{Shruti Priya, Samarendra Sahoo}
\date{Source: arXiv:2609.01372v1 (CC BY 4.0)}
\begin{document}
\maketitle

\noindent\emph{Source and licence.} Integrally closed ideals with $e_{2}(I)=e_{1}(I)-e_{0}(I)+\lambda(A/I)$. Version arXiv:2609.01372v1 (CC BY 4.0), distributed by arXiv under the Creative Commons Attribution 4.0 International licence, http://creativecommons.org/licenses/by/4.0/, as recorded in the arXiv licence metadata for this exact version and shown on the abstract page https://arxiv.org/abs/2609.01372v1. Full licence notice: \texttt{licenses/arxiv-2609.01372-CC-BY-4.0.txt}.\par\smallskip

\noindent\textit{Editorial scope.} Counted unit: the article's proposition e3bd (source0.tex lines 354--361). The article's complete original proof of it is delivered with it (source0.tex lines 363--402, one \texttt{proof} environment of 40 lines). No mathematics is written, repaired or re-derived: the statement and the proof are the article's own lines, in the article's own notation, and no retained line is substituted or re-flowed.  Retained uncounted as the source-local context the proof needs: lines 324--326.  Nothing was omitted from any retained span, so the package carries no omission record.  The retained lines cite 2 works, whose entries are transcribed verbatim from the article's own bibliography source in the e-print and delivered with short editorial labels recorded in the item's reference mapping.  No retained line contains a figure environment, an image-inclusion command or drawing code, so no image is delivered and the package declares no asset dependency.  Editorial formatting only, in addition: an AMS \texttt{amsart} preamble that restates the article's own declarations and macros used by the retained lines, namely the environments \texttt{lemma}, \texttt{proposition} on separate counters, together with the standard packages those lines need.  The retained environments print in this document as Lemma 1, Proposition 1 in order of appearance, whereas the article prints the same items with the numbering of its own full document.  The complete native source of the article as distributed in the arXiv e-print for this exact version is bundled unchanged as \texttt{sources/209105/source0.tex}.
\begin{lemma} \label{lemma4.4}
    Let $(A,\m)$ be a two-dimensional Cohen-Macaulay local ring and let $I$ be an $\m$-primary integrally closed ideal. If $e_{2}(I)=e_{1}(I)-e_{0}(I)+\lambda(A/I)$  then $\depth G(I)=0$ or $2$.
\end{lemma}

\begin{proposition}\label{e3bd}
Let $(A,\m) $ be a Cohen-Macaulay local ring of dimension $d \geq 2$ and let $I$ be an $\m$-primary integrally closed ideal. If $e_2(I)=e_1(I)-e_0(I)+\lambda(A/I)\neq 0$, then the following hold:
\begin{enumerate} [\normalfont(i)]
    \item \label{e3bd(i)} For $d=2,$ $e_3(I)\leq 0.$ Further, if  $e_3(I)=0$ then $G(I)$ is Cohen-Macaulay.

    \item \label{e3bd(ii)} For $d\geq 3,$ $e_3(I)\leq 0.$ Further, if $d=3$ and $e_3(I)=0$ then $\wt{G}(I)$ is Cohen-Macaulay. % In Particular, the Ratliff-Rush filtration with respect to $I$ behaves well modulo a superficial sequence of length $d-1$.
\end{enumerate}
\end{proposition}

\begin{proof} 


   (\ref{e3bd(i)}) Let $J$ be a minimal reduction of $I.$ Set $\wt{\nu_j}(I)=\lambda(\wt{I^{j+1}}/J\wt{I^j})$ for all $j\geq 0.$ From \cite[Equation 3.1]{rossi2010hilbert}, we have $e_1(I)=\sum_{j\geq 0}\wt{\nu_j}(I)$ and $e_2(I)=\sum_{j\geq 1}j\wt{\nu_j}(I).$ From the hypothesis, we obtain $\wt{\nu_j}(I)=0$ for all $j\geq 2.$ Thus, from \cite[Theorem 2.5$(c)$]{rossi2010hilbert}, we have
\begin{equation}\label{al}
    \wt{e}_{3}(I)=\displaystyle \sum _{j \geq 2}\binom{j}{2}\wt{\nu_j}(I)=0.
\end{equation}
 Furthermore, from \cite[1.5$(b)$]{Part2}, we also have
\begin{equation} \label{e3}
    e_{3}(I)=\wt{e}_{3}(I)-\sum_{j\geq 0}\lambda\left(\frac{\widetilde{I^{j+1}}}{I^{j+1}}\right)=-\sum_{j\geq 0}\lambda\left(\frac{\widetilde{I^{j+1}}}{I^{j+1}}\right)\leq 0.
\end{equation}
    Suppose that equality holds in (\ref{e3}), then $\widetilde{I^{j}}=I^{j}$ for all $j \geq 1.$ This implies $\depth G(I) >0.$ Then the result follows from Lemma \ref{lemma4.4}.


\noindent
(\ref{e3bd(ii)}) We proceed by induction on the dimension $d$. Suppose $d=3.$ We choose $x \in I\setminus I^2$ a superficial element for $I,$ such that $I/(x)$ is integrally closed. Set ${A_1}=A/(x)$ and ${I_1}=I/(x)$. From \cite[Proposition 1.2]{rossi2010hilbert}, we have $e_i(I)=e_i(I_1)$ for $i=0,1,2$ and
   \begin{align} \nonumber \label{123}
       e_{3}(I)&=e_{3}({I_1})+\sum_{j\geq 0}\lambda\left((I^{j+1}:x)/I^{j}\right)\\
       &=\widetilde{e}_{3}({I_1})-\sum_{j\geq 0}\lambda\left(\widetilde{{I^{j+1}_1}}/{I^{j+1}_1}\right)+\sum_{j\geq 0}\lambda\left((I^{j+1}:x)/I^{j}\right) \text{ (from \cite[1.5$(b)$]{Part2})}.
   \end{align}
    
 For all $j\geq 0$, we have the following exact sequence: 
\begin{equation}\label{behave}
    0\longrightarrow (I^{j+1}:x)/I^j\longrightarrow  \wt{I^j}/I^j \longrightarrow  \wt{I^{j+1}}/I^{j+1} \xlongrightarrow{g_{j+1}} \wt{{I^{j+1}_1}}/{I^{j+1}_1} \longrightarrow \mathrm{Coker}(g_{j+1}) \longrightarrow 0
\end{equation}
Taking lengths and summing over all $j \geq 0$,   we get:
\begin{equation*}
    \sum_{j\geq 0}\lambda\left((I^{j+1}:x)/I^j\right)+\sum_{j\geq 0}\lambda(\wt{I^{j+1}}/I^{j+1})-\sum_{j\geq 0}\lambda(\wt{I^j}/I^j)-\sum_{j\geq 0}\lambda\left(\wt{{I^{j+1}_1}}/{I^{j+1}_1}\right) \leq 0.
\end{equation*}
This implies that $\displaystyle\sum_{j\geq 0}\lambda\left((I^{j+1}:x)/I^j\right)-\sum_{j\geq 0}\lambda\left(\wt{{I^{j+1}_1}}/{I^{j+1}_1}\right)\leq 0.$ Note that the inequality is strict if the map $g_{j}$ is not surjective for some $j\geq 1.$
 Now, from equations \eqref{al} and (\ref{123}), we obtain 
    \begin{equation} \label{e30}
        e_3(I)=\widetilde{e}_{3}({I_1})-\sum_{j\geq 0}\lambda\left(\widetilde{{I^{j+1}_1}}/{I^{j+1}_1}\right)+\sum_{j\geq 0}\lambda\left((I^{j+1}:x)/I^{j}\right)\leq \widetilde{e}_{3}({I_1})=0.
    \end{equation}
    
    Suppose equality holds in (\ref{e30}), then the map $g_j$ is surjective for all $j\geq 1.$ Therefore, from \cite[Definition 4.4]{Part2}, the Ratliff-Rush filtration of $I$ behaves well modulo $(x),$ this implies $\depth \widetilde{G}(I) \geq 2.$ Furthermore, from \cite[Remark 4.5]{Part2}, $\widetilde{G}(I)/x\widetilde{G}(I)=\widetilde{G}({I_1}).$ Since $\widetilde{G}(I_1)$ is Cohen-Macaulay by \cite[Proposition 3.4$(iv)$]{Part2}, it follows from Sally's descent that $\widetilde{G}(I)$ is Cohen-Macaulay. 

    Assume that the statement is true for $d-1.$ Suppose $d\geq 4.$ Let $x \in I\backslash I^{2}$ be a superficial element for $I$ such that $I/(x)$ is integrally closed.
    From \cite[Proposition 1.2]{rossi2010hilbert}, $e_{3}(I)=e_{3}(I_1)$. By induction hypothesis, $e_3(I)\leq 0.$ %If equality holds, then $e_3(I_1)=0$. and from \eqref{behave} and \eqref{e30}, we obtain $g_j=0$ for all $j\geq 1.$ The result now follows from the induction hypothesis and the same argument as in the case $d=3.$ 
\end{proof}

\begin{thebibliography}{99}
\bibitem[Part2]{Part2} \bysame, \emph{Ratliff-Rush filtration, regularity and depth of higher associated graded modules. Part II}, J. Pure Appl. Algebra \textbf{221} (2017), no. 3, 611–631. %
\bibitem[rossi2]{rossi2010hilbert} M. E. Rossi and G. Valla, \emph{Hilbert functions of filtered modules}, vol. 9, Springer Science \& Business Media, 2010.
\end{thebibliography}
\end{document}
